% Short answer key for Listas/lista-r8b-executive-compensation.md (15 exercises).
% Compile: pdflatex twice; check: pdffonts (Type 1 only, uni = yes).
% Every number is asserted in lista-r8b-executive-compensation-solutions_check.py.
\documentclass[11pt,a4paper]{article}
\usepackage[utf8]{inputenc}
\usepackage[T1]{fontenc}
\usepackage{lmodern}   % Type 1 vector fonts (replaces PK bitmaps)
\usepackage{cmap}      % ToUnicode CMap -> copyable text and accents
\usepackage[brazilian,english]{babel}
\usepackage{amsmath,amssymb}
\usepackage[a4paper,margin=2.2cm]{geometry}
\usepackage{enumitem}
\usepackage{fancyhdr}
\usepackage{microtype}

\pagestyle{fancy}\fancyhf{}
\lhead{\footnotesize Finan\c{c}as II -- Corporativas (MPE Insper)}
\rhead{\footnotesize List R8b --- Answer key}
\cfoot{\footnotesize\thepage}
\renewcommand{\headrulewidth}{0.3pt}
\setlength{\headheight}{14pt}
\setlist{topsep=2pt,itemsep=1pt}
\setlength{\parindent}{0pt}
\newcommand{\dpp}{\Delta p}

\usepackage[hyphens]{url}
\begin{document}
\begin{center}{\Large\bfseries Exercise list R8b --- Executive compensation and managerial incentives}\\[4pt]{\large Answer key (full derivations in \url{lista-r8b-executive-compensation-solutions.pdf})}\end{center}

Amounts in \$ millions. Numbers marked ``illustrative'' are parameter points chosen in the full solutions.

\section*{1. Is CEO pay too high? (Ross CQ 1.16)}
\$550m over five years is \$110m a year; \$54.4m is 1\% of a \$5.44bn value gain. Not necessarily excessive:
scarce talent applied to a very large firm, as with sports stars. Excessive if not set at arm's length
(boards influenced by the CEO) or if it pays for industry luck (the oil price) rather than value created
relative to peers. Jensen--Murphy: \$3 per \$1{,}000 --- the problem is weak sensitivity, not level.

\section*{2. GreenFrame acquisition (Berk Ch.~22, Problem 3)}
CEO's loss $0.03\times50=1.5$; net $5-1.5=+3.5$: \textbf{better off}. He would refuse only with a stake above
$5/50=10\%$, or if the deal destroyed more than $5/0.03=166.7$.

\section*{3. Why managers want growth (G\&T 18.1)}
Expand iff $\Delta W+b>\alpha|\mathrm{NPV}|$: pay and prestige rise with size, perquisites, job security from
diversification (43\% of CEOs keep their jobs after bankruptcy), careers for staff, entrenchment; free cash
flow finances it (RJR: \$148m on \$2.8bn, 5.3\%). Remedies: value-based pay, ownership, debt.

\section*{4. Braintree bonus (Higgins 1.2)}
ROE $=100/5{,}000=2\%$, ROA $=1.43\%$; equity charge at 10\% $=500$; economic profit $100-500=-400$. The
profit came from the boom (luck). \textbf{Reject}; tie any bonus to economic profit relative to peers.

\section*{5. Stock pay and market-to-book (G\&T 18.7)}
(1) Incentives: stock-based pay raises value (Mehran; +20\% at adoption). (2) Firm type: high-M/B growth
firms have intangibles, discretion and poor accounting measures, so stock pay is the efficient contract.
(Also signalling: options at \$20 are worth 10 to Jack (value 30), 0 to Peggy (value 15).)

\section*{6. Chicago Bulls coach (G\&T 18.9)}
IC: bonus $\ge B/\Delta_s$. Illustrative: wins $0.6/0.4\Rightarrow b=5$, expected pay 3; experts $0.9/0.3\Rightarrow
b=1.667$, expected pay 1.5. (a) Wins: aligned, objective, cheap; but luck/roster risk, short-termism.
(b) Decisions: filters luck, rewards long-term choices; but costly, subjective, non-verifiable,
conformist. Best: wins relative to a benchmark plus an evaluation component.

\section*{7. Chrysler vs Chevron (G\&T 18.6)}
Chevron: filter the oil price (indexed/peer-relative pay; illustrative $R^2=0.7$ lets the slope rise
$1/\sqrt{0.3}=1.83\times$), charge for capital to curb free-cash-flow overinvestment. Chrysler: high
stock-based sensitivity (value depends on the CEO's strategy), limited explicit relative performance
(aggressive competition among three automakers), watch risk-taking near distress.

\section*{8. High Tech vs Super Tech (G\&T 18.3)}
(1) Jacobs $25.5/75=34\%$: loses majority, bears less of perks; unit managers paid in diluted High Tech
stock. (2) Keeps 51\% of High Tech; look-through $0.51\times0.55=28.05\%$ of Super Tech; unit managers paid
in Super Tech stock (stronger incentives, market monitoring); but Jacobs gains $0.51-0.2805=0.2295$ per
\$1 shifted from Super Tech to High Tech.

\section*{9. Cybertex recap (G\&T 18.8)}
$V=100$: managers' $0.5$ of equity becomes $0.5/20=2.5\%$ (from 1\%). (i) Fewer pet projects (2.5 times
the cost borne; debt drains free cash). (ii) Ambiguous: levered equity rewards risk, but distress risk and
undiversified human capital make managers conservative. (iii) Fewer long-term projects (overhang,
cash commitments, distress fear).

\section*{10. Stock options in \S3.1.1 (Amaro 3.1.5 \#3)}
FOC $-\alpha U_1+U_2=0$; $dF/d\alpha=U_1/g_F<0$ (no wealth effects). Options raise the marginal stake to
$\alpha+n\Delta$, like a higher $\alpha$: illustrative $U=\alpha(100-F)+\sqrt F$, $F^*=(0.5/\alpha)^2$: $V$ from 75
($\alpha=0.10$) to 88.89 ($\alpha_{\text{eff}}=0.15$). Add: risk-taking and bondholder transfers, dilution,
falling delta, managerial risk aversion, price noise, manipulation. Matches the book (p.~256).

\section*{11. Condition (3.4) (Amaro 3.1.5 \#4)}
$U'[w-f]=\lambda G'[f]$: equal MRS across states (Pareto-efficient risk sharing). Differentiating:
$\dfrac{df}{d\theta}=\dfrac{R^P}{R^P+R^A}\dfrac{dw}{d\theta}$; constant fee if shareholders are risk neutral, full
payoff minus a rent if the manager is. Matches the book (pp.~256--257).

\section*{12. CARA $\Rightarrow$ linear fee (Amaro 3.1.5 \#6)}
$e^{-\delta(w-f)}=\lambda e^{-\alpha f}\Rightarrow f=\dfrac{\delta}{\alpha+\delta}w+\dfrac{\ln\lambda}{\alpha+\delta}$. Similarity also holds:
$U=(\alpha\lambda/\delta)G$. Matches the book (p.~257).

\section*{13. Congruence and power (Tirole 3.10)}
High-powered: NPV$^1=p_HR-I+(1-x)B$, $\mathcal P^1=p_H(R-B/\dpp)$. Low-powered ($R_b=0$):
$\mathcal P^2=p_HR-x\dpp R$, NPV gap $x(\dpp R-B)>0$. Low-powered optimal only if (1) fails, (2) holds,
which needs $x<\hat x=p_HB/[(\dpp)^2R]$; for $x\ge\hat x$ any funded scheme is high-powered. Illustrative
($\hat x=0.45$): $x=0.2$, $I-A=7.6$: $7.2<7.6\le8.2$, low-powered. (ii) The congruent type can take the
dissonant scheme, so rents are at least $p_HB/\dpp$ (or 0): menus cannot beat simple schemes. Matches
the book (p.~575).

\section*{14. Second best is not Pareto-efficient (Amaro 3.1.5 \#5)}
Risk-neutral shareholders: Pareto efficiency = constant fee = no effort. Incentives need $f$ to vary, so
$U'/G'$ varies ($\psi>0$). Illustrative ($G=\sqrt f$): $f_H=30.25$, $f_L=9$, pay 26 vs 25 first best;
$1/G'=11,6$; $\psi=2$, $\lambda=10$. Risk-neutral manager: sell him the firm ($K=29$), $\psi=0$,
Pareto-efficient. Matches the book (p.~257).

\section*{15. Golden parachutes (Tirole 7.9)}
$R_b^r=B/\dpp$, $T^r=0$; only the $q$-type's rent $\Pi\ge\max\{z^r(q_H-\tau)b,\,z^qq_Hb\}$ matters.
$q_HR>L$: $T^{q*}=0$ (continuation is the cheaper currency, $q_HR_b+dT/dz^q=q_HR-L>0$); under (7.10)
funding fails if $(1,z^*)$ violates (NM). $q_HR<L$: $z^{q*}=0$, $T^{q*}=z^r(q_H-\tau)B/\dpp$; illustrative
$T^*=0.6$, NPV 1.5 vs 1.333 without the parachute. Matches the book (pp.~604--605).

\end{document}
